\subsection{Construcción correlacionada y orden de los lectores}
\label{gen:orden-constructor-correlacionado}

El antecedente común de las tres regiones es la emisión aritmética con
su registro de trayectoria. APP proporciona las evaluaciones suma y
producto; TRIT aporta la firma orientada; la transición
\eqref{eq:actualizacion} compone selección de modo, transporte cardinal y
actualización de memoria. El censo de esa emisión produce el catálogo y
sus fibras. Los predicados de multiplicidad, soporte diagonal y orientación
seleccionan los representantes de \eqref{eq:palabras-regionales}.
La correlación regional conserva esta procedencia conjunta, junto con las
incidencias y las restricciones entre sus registros.

Los transportes \(L_0,L_1\) actúan sobre los representantes seleccionados
en la coordenatización fila declarada. El calendario y la concatenación distinguen
la longitud del prefijo, la fase visible y la memoria acumulada. La
coordenada profunda y sus cocientes se explicitan a continuación mediante
\eqref{memprof:division} y \eqref{memprof:inversa}. Las restricciones de
historias se conservan en el sentido de \ref{hist:seccion}.
La selección orientada concreta de \(R_{36}\) tiene su demostración
incidencial en \ref{mg01:frontera}, después de la construcción de \(A_W\);
esa residencia conserva los antecedentes utilizados por su prueba.

Cada lector regional añade una operación tipada al registro que recibe.
En la clausura, el lector simétrico de la pentafibra fija la pendiente
primitiva y la ley de composición determina el compensador de
\eqref{eq:lector-pi}. En la propagación, la composición formal, la unidad
y el generador orientado determinan la recurrencia
\eqref{eq:propagacion}. En la autoescala, el calendario de modos produce
la incidencia \eqref{eq:fav}, cuyo rayo positivo fija la coordenada
normalizada. Las demostraciones siguientes conservan las normalizaciones
y los dominios de estas tres operaciones.

Las cotas racionales y la separación de cilindros evalúan después los
caracteres obtenidos. La emisión posicional efectiva y su compatibilidad
están demostradas en \ref{act21:thm:df-emision-correcta} y
\ref{act21:prop:df-prefijos-compatibles}. El registro dodecafásico conserva
su construcción propia en \ref{act21:chap:despliegue-selector-k};
sus coordenadas y las tres lecturas regionales intervienen posteriormente
en la clausura entera de \ref{sec:alpha}. De este modo, selección de
regiones, transporte, construcción de caracteres y evaluación de prefijos
mantienen sus dominios y su orden, dentro de la misma genealogía
APP--TRIT--TPK.
